\documentclass[11pt]{article}
\usepackage[a4paper,margin=1in]{geometry}
\usepackage{amsmath,amssymb,amsthm,bm}
\usepackage{booktabs}
\usepackage{enumitem}
\usepackage{hyperref}

\title{Data Mining Task: Problem Formulation and Prediction Objective}
\author{}
\date{}

\begin{document}
\maketitle

\section{Problem Setup}
Given papers and extracted phrases over years $t\in[t_{\min},t_{\max}]$, define yearly transaction sets
\begin{equation}
\mathcal{T}_t = \{\tau_{t,1},\dots,\tau_{t,N_t}\},
\end{equation}
where each transaction is a set of phrase IDs appearing in one paper.

Goal:
\begin{enumerate}[leftmargin=1.5em]
  \item mine frequent/strong itemsets (pairs and triples) per year,
  \item build itemset time series,
  \item compute trend residuals against rolling baselines,
  \item predict next-year trend score and rank itemsets.
\end{enumerate}

\section{Frequent Pair Mining (Per Year)}
For year $t$, support count and fraction for pair $s=\{i,j\}$:
\begin{align}
\mathrm{sup\_count}_t(s) &= \left|\{\tau\in\mathcal{T}_t: i\in\tau,\;j\in\tau\}\right|,\\
\mathrm{sup}_t(s) &= \frac{\mathrm{sup\_count}_t(s)}{N_t}.
\end{align}

A pair is frequent if
\begin{equation}
\mathrm{sup\_count}_t(s) \ge \max(\texttt{min\_count},\lceil \texttt{minsup\_frac}\cdot N_t\rceil).
\end{equation}

Lift of a pair:
\begin{equation}
\mathrm{lift}_t(i,j)=\frac{\mathrm{sup}_t(\{i,j\})}{\mathrm{sup}_t(\{i\})\,\mathrm{sup}_t(\{j\})+\varepsilon}.
\end{equation}
Strong pairs satisfy $\mathrm{lift}_t(i,j)\ge \texttt{min\_lift}$.

\section{Triple Mining (Per Year)}
For triple $s=\{i,j,k\}$, count and support are computed analogously.
Candidate pruning uses pairwise strong-pair constraints; keep triples that satisfy minimum support and lift constraints.

\section{Itemset Time Series Construction}
Construct an itemset universe $\mathcal{S}$ from strong pairs/triples across years with activity constraints (e.g., minimum active years and recent activity).

For each $s\in\mathcal{S}$ and year $t$, define support series:
\begin{equation}
\mathrm{sup}_t(s).
\end{equation}
Apply log transform:
\begin{equation}
z_t(s)=\log(1+\mathrm{sup}_t(s)).
\end{equation}

\section{Rolling Baseline and Residual Trend}
For each itemset $s$, fit a rolling baseline on the previous $L$ years (method = drift or AR(1)).
Predicted baseline:
\begin{equation}
\widehat{z}_t(s)=\text{Baseline}\big(z_{t-L}(s),\dots,z_{t-1}(s)\big).
\end{equation}
Residual trend signal:
\begin{equation}
r_t(s)=z_t(s)-\widehat{z}_t(s).
\end{equation}
A larger positive residual indicates unexpectedly strong growth relative to baseline behavior.

\section{Next-Year Prediction Objective}
For target year $t_0$ (e.g., 2025) and next-year forecast $t_0+1$ (e.g., 2026), predict residual via recent-window average:
\begin{equation}
\widehat{r}_{t_0+1}(s)=\frac{1}{|\mathcal{W}|}\sum_{\tau\in\mathcal{W}} r_{\tau}(s),
\end{equation}
where $\mathcal{W}$ is a residual window (e.g., recent years up to $t_0$).

Define recent support statistic:
\begin{equation}
\mathrm{recent\_support}(s)=\frac{1}{|\mathcal{W}|}\sum_{\tau\in\mathcal{W}}\mathrm{sup}_{\tau}(s).
\end{equation}

Ranking scores:
\begin{align}
\mathrm{score}(s) &= \widehat{r}_{t_0+1}(s),\\
\mathrm{score\_weighted}(s) &= \widehat{r}_{t_0+1}(s)\cdot\big(1+\mathrm{recent\_support}(s)\big).
\end{align}

Itemsets are ranked descending by score (or weighted score) to produce top predicted trends.

\section{Outputs}
The task returns:
\begin{itemize}[leftmargin=1.5em]
  \item top frequent pairs by support,
  \item top pairs by lift,
  \item top strong triples,
  \item predicted next-year trending itemsets with residual-based ranking,
  \item optional keyword-focused support/residual time series.
\end{itemize}

\end{document}
